\documentclass[12pt]{article}
\usepackage[utf8]{inputenc}
\usepackage{geometry}
\geometry{letterpaper, margin=1in}
\usepackage{enumitem}
\usepackage{tabularx}
\usepackage{hyperref}

\begin{document}

% --- 1. Cover Page ---
\begin{titlepage}
    \centering
    \vspace*{2in}
    {\Huge \textbf{Software Requirement Specification (SRS)}}\\[1cm]
    {\Large \textbf{EagleNav: Campus Navigation System}}\\[2cm]
    {\large \textbf{Team Name:} Moh's Eagles}\\[0.5cm]
    {\large \textbf{Date:} May 15, 2026}\\[2cm]
    \vfill
\end{titlepage}

% --- 2. Table of Contents ---
\tableofcontents
\newpage

% --- 3. Version Description ---
\section*{Version Description}
\addcontentsline{toc}{section}{Version Description}
\begin{table}[h]
    \centering
    \begin{tabularx}{\textwidth}{|c|X|c|}
        \hline
        \textbf{Version Number} & \textbf{Description of what was added} & \textbf{Date Added} \\ \hline
        0.0 & Initial SRS creation. & Feb 13, 2026\\ \hline
        1.0 & Outlining of core interface and routing requirements. & Apr 27, 2026 \\ \hline
        2.0 & Added functional requirements for Events Subsystem notifications. & May 4, 2026 \\ \hline
        3.0 & Added hardware constraints for Computer Vision (Camera) integration. & May 11, 2026 \\ \hline
        4.0 & Finalized ethical considerations and Landmark Service audio requirements. & May 15, 2026 \\ \hline
    \end{tabularx}
\end{table}

\newpage

% --- 4. Sections ---

\section{Introduction}

\subsection{Purpose of the document}
This Software Requirement Specification (SRS) defines the functional, non-functional, and ethical requirements for the EagleNav navigation system. It establishes the baseline expectations for the application's external interfaces across all project phases.

\subsection{Intended audience}
This document is intended for software developers, UI/UX designers, and quality assurance testers to ensure the system is built to meet the accessibility, safety, and privacy standards required for a campus utility application.

\subsection{Overview of the software}
EagleNav provides high-precision pedestrian navigation. The mobile client communicates with an external cloud routing engine to deliver visual and logical guidance to accessible entrances. It utilizes device sensors (Camera, GPS, Magnetometer) to provide real-time spatial awareness, obstacle detection, and localized campus event alerts.

\section{External Interface Requirements}

\subsection{User Interface}
\begin{itemize}[noitemsep]
    \item The software must present a standard interactive map with a polyline path overlay.
    \item The UI must be fully compatible with native OS screen reading software (VoiceOver for iOS, TalkBack for Android).
    \item The device must utilize specific haptic vibration patterns to indicate turns or obstacle warnings generated by the AR camera feed.
    \item The Text-to-Speech (TTS) manager must enforce priority queuing so urgent obstacle warnings interrupt standard navigation directions.
\end{itemize}

\subsection{Software Interfaces}
\begin{itemize}[noitemsep]
    \item \textbf{Routing API:} The software will communicate with a Dockerized Valhalla routing server via HTTP POST requests containing origin, destination, and pedestrian costing options (e.g., "avoid stairs").
    \item \textbf{Event Scraper:} The software must fetch HTML from the CSULA events page, parse it into structured JSON objects, and interface with the OS notification scheduler for local alerts.
\end{itemize}

\section{Legal and Ethical Considerations}

\subsection{Data storage and privacy considerations}
EagleNav operates with a strict privacy-first approach. The system does not require user accounts or logins. No Personally Identifiable Information (PII), student records, search history, or location logs will be transmitted to or stored on the cloud routing server. All routing computations are anonymous. 

\subsection{Any legal or ethical issues}
\begin{itemize}[noitemsep]
    \item \textbf{User Consent:} The software must explicitly request and secure the user's permission to access device location services and the device camera upon first launch.
    \item \textbf{ADA Compliance:} The software is built under the guidelines of the Americans with Disabilities Act to ensure equal access to campus facilities.
    \item \textbf{Safety Liability:} The software must present a legal disclaimer indicating that it is an assistive tool and not a replacement for physical mobility aids (e.g., canes, guide dogs). Users are ultimately responsible for their physical safety while traversing the campus.
\end{itemize}

\section{Glossary}
\begin{itemize}[noitemsep]
    \item \textbf{ADA:} Americans with Disabilities Act.
    \item \textbf{JSON:} JavaScript Object Notation.
    \item \textbf{PII:} Personally Identifiable Information.
    \item \textbf{SRS:} Software Requirement Specification.
\end{itemize}

\end{document}